\documentclass[../preamble]{subfiles}
\begin{document}

\section{Braided Monoidal Categories}
\label{sec:braided-monoidal}
\concept{braided-monoidal}{Braided Monoidal Category}
\depends{monoidal-category}
\implements{haskell}{Cat.BraidedMonoidal}
\implements{agda}{Cat.BraidedMonoidal}
\implements{lisp}{braided-monoidal}

A \newterm{braided monoidal category} equips a monoidal category
with a coherent notion of ``swapping'' the two factors of a tensor product
(\nlabref{braided+monoidal+category};
Joyal--Street, \emph{Braided tensor categories}, Advances in Mathematics 102, 1993).

\begin{definition}[Braided monoidal category]
  \label{def:braided-monoidal}
  A \newterm{braided monoidal category}
  is a monoidal category
  $(\Category{C}, \otimes, I, \alpha, \lambda, \rho)$
  (Definition~\ref{def:monoidal-category})
  equipped with a natural isomorphism,
  the \newterm{braiding},
  with components
  \begin{equation}
    \label{eq:braiding}
    \newmath{\sigma_{A,B}
      \colon A \otimes B
      \xrightarrow{\;\sim\;} B \otimes A}
  \end{equation}
  natural in $A$ and $B$,
  subject to the two hexagon axioms
  (Definition~\ref{def:hexagon-1} and Definition~\ref{def:hexagon-2}).
\end{definition}

\subsection{Hexagon Axioms}

\begin{definition}[First hexagon axiom]
  \label{def:hexagon-1}
  \axiom{braided-monoidal}{hexagon-1}
  The \newterm{first hexagon axiom} requires that
  for all objects $A, B, C \in \Ob(\Category{C})$,
  the following diagram commutes:
  \begin{equation}
    \label{eq:hexagon-1}
    \begin{tikzcd}[column sep=large]
      (A \otimes B) \otimes C
        \arrow[r, "\alpha_{A,B,C}"]
        \arrow[d, "\sigma_{A,B} \otimes \id_C"']
      & A \otimes (B \otimes C)
        \arrow[r, "\sigma_{A, B \otimes C}"]
      & (B \otimes C) \otimes A
        \arrow[d, "\alpha_{B,C,A}"]
      \\
      (B \otimes A) \otimes C
        \arrow[r, "\alpha_{B,A,C}"']
      & B \otimes (A \otimes C)
        \arrow[r, "\id_B \otimes \sigma_{A,C}"']
      & B \otimes (C \otimes A)
    \end{tikzcd}
  \end{equation}
  That is,
  the two composite paths from
  $(A \otimes B) \otimes C$
  to
  $B \otimes (C \otimes A)$
  are equal.
\end{definition}

\begin{definition}[Second hexagon axiom]
  \label{def:hexagon-2}
  \axiom{braided-monoidal}{hexagon-2}
  The \newterm{second hexagon axiom} requires that
  for all objects $A, B, C \in \Ob(\Category{C})$,
  the following diagram commutes:
  \begin{equation}
    \label{eq:hexagon-2}
    \begin{tikzcd}[column sep=large]
      A \otimes (B \otimes C)
        \arrow[r, "\alpha^{-1}_{A,B,C}"]
        \arrow[d, "\id_A \otimes \sigma_{B,C}"']
      & (A \otimes B) \otimes C
        \arrow[r, "\sigma_{A \otimes B, C}"]
      & C \otimes (A \otimes B)
        \arrow[d, "\alpha^{-1}_{C,A,B}"]
      \\
      A \otimes (C \otimes B)
        \arrow[r, "\alpha^{-1}_{A,C,B}"']
      & (A \otimes C) \otimes B
        \arrow[r, "\sigma_{A,C} \otimes \id_B"']
      & (C \otimes A) \otimes B
    \end{tikzcd}
  \end{equation}
  That is,
  the two composite paths from
  $A \otimes (B \otimes C)$
  to
  $(C \otimes A) \otimes B$
  are equal.
\end{definition}

\subsection{Remarks}

\begin{remark}[Naturality of the braiding]
  \label{rem:braiding-naturality}
  The naturality of $\sigma$ means that
  for all morphisms
  $f \colon A \lto A'$ and $g \colon B \lto B'$,
  the following square commutes:
  \begin{equation}
    \label{eq:braiding-naturality}
    \begin{tikzcd}
      A \otimes B
        \arrow[r, "\sigma_{A,B}"]
        \arrow[d, "f \otimes g"']
      & B \otimes A
        \arrow[d, "g \otimes f"]
      \\
      A' \otimes B'
        \arrow[r, "\sigma_{A',B'}"']
      & B' \otimes A'
    \end{tikzcd}
  \end{equation}
\end{remark}

\begin{remark}[Naming convention]
  \label{rem:hexagon-naming}
  The hexagon axioms are so named because,
  when expanded using explicit associators,
  each diagram has six edges forming a hexagonal shape.
  Our rectangular presentation is equivalent
  and more compact.
\end{remark}

\begin{remark}[Yang--Baxter equation]
  \label{rem:yang-baxter}
  In a braided monoidal category,
  the braiding satisfies the
  \newterm{Yang--Baxter equation}:
  \begin{equation}
    \label{eq:yang-baxter}
    (\id_B \otimes \sigma_{A,C})
    \circ (\sigma_{A,B} \otimes \id_C)
    \circ (\id_A \otimes \sigma_{B,C})
    =
    (\sigma_{B,C} \otimes \id_A)
    \circ (\id_B \otimes \sigma_{A,C})
    \circ (\sigma_{A,B} \otimes \id_C)
  \end{equation}
  as morphisms
  $A \otimes B \otimes C \lto C \otimes B \otimes A$
  (where we suppress associators by Mac Lane's coherence theorem).
  This follows from the hexagon axioms alone.
\end{remark}

\end{document}
